\documentclass[11pt]{article}
\usepackage[margin=1in]{geometry}
\usepackage{amsmath,amssymb,amsthm,microtype}
\usepackage{xcolor}
\usepackage[colorlinks=true,allcolors=blue!45!black]{hyperref}
\hypersetup{pdftitle={Separated intervals with reciprocal sum one}}
\newtheorem{theorem}{Theorem}
\newtheorem{lemma}{Lemma}
\newcommand{\Z}{\mathbb Z}
\newcommand{\Q}{\mathbb Q}
\DeclareMathOperator{\lcm}{lcm}
\setlength{\emergencystretch}{2em}
\allowdisplaybreaks
\title{Separated intervals with reciprocal sum one}
\author{Stijn Cambie \and Johan Land \and Yuren Tang}\date{}
\makeatletter
\renewcommand{\maketitle}{%
 \begin{center}{\LARGE\@title\par}\end{center}\vspace{.35\baselineskip}}
\makeatother
\begin{document}
\maketitle
\begin{abstract}
For every sufficiently large \(k\), we construct \(k\) nonadjacent
integer intervals, each of length at least two, whose reciprocal sum
is \(1\). The proof uses elementary prime estimates, the marriage
theorem, and a short roots-of-unity argument. 
% Its main observation is
% that residues and the number of intervals can be controlled together.
\end{abstract}


\section{Introduction}


\textcolor{red}{real intro to be written}

Call a finite family of integer intervals in \(\{2,3,\ldots\}\)
a \emph{configuration} if each interval has length at least two
and successive intervals have at least one unused integer between them.
Write \(w(F)\) for its reciprocal sum and \(\nu(F)\) for its number of
intervals. We prove the following form of Erd\H{o}s Problem~289
\cite{EG}.

\begin{theorem}\label{main}
Every sufficiently large integer \(k\) admits a configuration \(F\)
with \(\nu(F)=k\) and \(w(F)=1\).
All intervals have length two, three or four; those longer than two
belong to a fixed finite set independent of \(k\).
\end{theorem}

Put \(D_n=\lcm(1,\ldots,n)\) and
\(G_n=D_n^{-1}\Z/\Z\subset\Q/\Z\).
Thus \(1/a\in G_n\) precisely when every prime power dividing \(a\)
is at most \(n\). The group changes only at prime powers:
if \(q=p^\alpha\), then \(G_q/G_{q-1}\) has order \(p\).
We use this organization, also present in \cite{Land,Tang}, to
realize whole sets \((\beta+G_n)\times I_n\) of attainable
\emph{residues and counts}: every residue in \(\beta+G_n\)
occurs at every integer count in \(I_n\).
All weights will be positive and below \(2\).
Once the fixed rational translate \(\beta\) lies in \(G_n\),
residue zero therefore gives weight exactly \(1\).

\section{Pairs supplying the successive quotients}

\begin{lemma}\label{pairs}
For every sufficiently large prime power \(q=p^\alpha\), there are
\(\gg q/\log q\) pairs \(\{qm,qm+\varepsilon\}\), with
\(\varepsilon=\pm1\), distinct \(1\le m<q\), \(p\nmid m\), and every
prime power dividing \(qm+\varepsilon\) smaller than \(q\).
Each pair lies in \([a-1,a+1]\), with distinct centers
\[
              a\equiv0\pmod8,\qquad cq^2/\log q\le a\le q^2
\]
for an absolute \(c>0\). Each center occurs at at most three stages \(q\).
\end{lemma}

\begin{proof}
For a sufficiently large fixed \(D\), Chebyshev's estimates
\(\pi(x)\asymp x/\log x\) \cite{HW} supply \(\gg q/\log q\) primes
\(b\in(q/D,q/16)\), excluding \(p\). Set \(d=8\) for odd \(p\)
and \(d=1\) for \(p=2\).
For \(r=db\), let \(t\in[1,q-1]\) be its inverse modulo \(q\).
Then
\[
 m_+=\frac{rt-1}{q},\qquad
 m_-=\frac{r(q-t)+1}{q},\qquad m_++m_-=r.
\]
Both coefficients lie in \((0,r)\), and at least one is coprime to
\(p\). Choose it. Its companion is \(dbT=qm+\varepsilon\), with
\(T=t\) or \(q-t\).

Only \(O(1)\) carrier primes \(b\) can fail the required
prime-power condition. Primes other than \(2,b\) contribute only
through \(T<q\). A power of \(b\) at least \(q\) requires
\(T=jb\), \(1\le j<D\), and hence
\(djb^2\equiv\varepsilon\pmod q\).
Each of these finitely many congruences has at most four roots:
if a solution exists, division by it reduces the question to
\(z^2\equiv1\pmod{p^\alpha}\), and one factors \((z-1)(z+1)\).
For \(d=8\), a bad power of two forces \(T\) to be a multiple
of the least power of two at least \(q/8\). There are at most
eight such \(T<q\), each determining at most one \(b<q\) per
sign through \(8bT\equiv\varepsilon\pmod q\).
For \(p=2\), both \(b,T\) are odd. Discard these exceptions.

Each \(m\) comes from at most four carriers: a carrier divides
\(qm-1\) or \(qm+1\), and each of these numbers, being below \(q^2\),
has at most two distinct prime factors greater than \(q/D\)
when \(q>D^3\).
After removing repetitions, retain the larger half of the
coefficients. There remain \(\gg q/\log q\), all
\(\gg q/\log q\).
Use \(a=qm+\varepsilon\) for odd \(p\), and \(a=qm\) for \(p=2\).
These are distinct multiples of \(8\) with the stated bounds.
Finally \(q\) is the largest prime-power divisor of \(qm\);
thus a center can occur only at the largest prime-power stages
of \(a-1,a,a+1\).
\end{proof}

\begin{lemma}\label{subsets}
For sufficiently large \(n\), the subset sums of any \(N\ge n^{3/4}\)
distinct units modulo \(n\) cover \(\Z/n\Z\).
\end{lemma}

\begin{proof}
For a unit set \(A\) and \(\zeta=e^{2\pi i/n}\), the number \(R(b)\)
of subsets summing to \(b\) satisfies
\[
 nR(b)=\sum_{t=0}^{n-1}\zeta^{-bt}\prod_{a\in A}(1+\zeta^{ta}).
\]
For \(0<t<n\), put \(g=\gcd(t,n)\). The residues \(ta\) are
nonzero multiples of \(g\), each with at most \(g\) preimages.
The two open end intervals of length \(N/4\) contain at most
\(N/(2g)\) such multiples altogether.
Hence at least \(N/2\) members satisfy
\(\|ta/n\|\ge N/(4n)\), where \(\|x\|\) is distance to the nearest integer.
Using \(|\cos\pi x|\le e^{-2\|x\|^2}\), we obtain
\[
 \left|\prod_{a\in A}(1+\zeta^{ta})\right|
       \le2^N e^{-N^3/(16n^2)}.
\]
The \(t=0\) term is \(2^N\), so
\(nR(b)\ge2^N(1-(n-1)e^{-n^{1/4}/16})>0\) for large \(n\).
\end{proof}

Every pair in Lemma~\ref{pairs} has weight in \(G_q\).
Modulo \(G_{q-1}\), its weight equals \(u/q\), where
\(u=m^{-1}\pmod q\): the companion lies in \(G_{q-1}\), and
\[
                 \frac uq-\frac1{qm}=\frac{um-1}{qm}
\]
has denominator dividing \(m\), hence dividing \(D_{q-1}\).
Consequently any \(s(q)=\lceil q^{3/4}\rceil\) of these pairs
have subset sums covering \(G_q/G_{q-1}\), by Lemma~\ref{subsets}.

\section{A finite starting family}

Let \(B\) be a sufficiently large power of two. All conditions
on its size below will be uniform in the later choices.
Write \(D_B=2^eM\), \(M\) odd, and choose \(N\ge e\) with
\(2^{N-3}\ge M\). Put \(K=M2^N\).
Reserve the following \(t=2(N-2)\) pairs:
\begin{equation}\label{chains}
 [2^i+1,2^i+2],\qquad[M2^i+1,M2^i+2],
                            \qquad 3\le i\le N.
\end{equation}
Their possible left extensions are \([h,h+2]\), with distinct
starts \(h\equiv0\pmod8\). Thus all are separated, and their
baseline weight \(\sigma\) is less than \(\frac12(1+1/M)\).
An extension adds \(1/h\). Multiplying by \(K\), the two chains
offer the increments \(M2^j\) and \(2^j\), \(0\le j\le N-3\).
They represent every integer \(0\le z<K/4\):
write \(z=Mu+v\), where \(u<2^{N-2}\) and
\(0\le v<M\le2^{N-3}\), and use the binary digits of \(u,v\).

Five single intervals provide the remaining coarse adjustment:
\[
\begin{array}{c|ccccc}
 h&0&15&28&43&55\\ \hline
 J_h &[5,6]&[4,6]&[2,3]&[2,4]&[2,5]\\
 w(J_h)&\multicolumn{5}{c}{11/30+h/60}.
\end{array}
\]
The successive gaps in \(\{0,15,28,43,55,60\}/60\) are at most
\(1/4\). Given \(x\in G_B\), represented in \([0,1)\), choose
the greatest \(h/60\le x\). The remainder belongs to
\([0,1/4)\cap K^{-1}\Z\), since \(60\mid D_B\mid K\).
The chains supply it. Each resulting configuration therefore has
\[
                  w=c_0+x,\qquad \nu=t+1,
       \qquad c_0=11/30+\sigma<9/10.
\]
The small interval ends at most at \(6\), while every possible
chain extension starts at least at \(8\), so separation is preserved.

Now take \(2s(B)\) pairs from Lemma~\ref{pairs} at stage \(B\),
avoiding all possible seed supports.
There are \(\gg B/\log B\) candidates. Only \(O(\log\log B)\)
can conflict with the two chains: in the stage band, whose endpoint
ratio is \(O(\log B)\), each geometric chain has that many terms,
uniformly in \(M,N\). Each term forbids \(O(1)\) centers.
The small intervals add \(O(1)\) exclusions.
Thus enough pairs remain, and their total weight is
\(O(\log B/B^{5/4})<1/20\).
Their weights lie in \(G_B\), so adding any fixed \(r\)-element
subset merely permutes the residues supplied by the seed.
For \(0\le r\le2s(B)\), this gives a finite family realizing
\begin{equation}\label{initial}
      (\beta+G_B)\times[u_B,u_B+2s(B)],
      \qquad \beta=c_0\pmod1,\quad u_B=t+1.
\end{equation}
All its weights are below \(19/10+1/20=39/20\).
Its supports and the rational translate \(\beta\) are now fixed.

\section{Enlarging the residue group and the count range}

We use one elementary observation. Suppose a family realizes
\((\beta+G)\times[u,v]\), and a new pool is separated internally
and from every old configuration. Its pair weights lie in
\(G'\supset G\); \(s\) pairs have subset sums covering \(G'/G\),
and there are \(m\) further pairs.
If \(v-u\ge s\), the unions realize
\begin{equation}\label{transfer}
                    (\beta+G')\times[u+s,v+m].
\end{equation}
Indeed, for a target count \(k\in[u+s,v+m]\), set
\(r=\max\{0,k-v\}\). Then \(0\le r\le m\) and
\(u+s\le k-r\le v\).
Choose any \(r\) further pairs \(A\).
For a target residue \(z\in\beta+G'\), choose a subset \(C\)
of the first \(s\) pairs with \(z-\beta-w(A)-w(C)\in G\).
The old family supplies residue \(z-w(A)-w(C)\) and count
\(k-r-|C|\in[u,v]\); adjoining \(A,C\) proves \eqref{transfer}.

At each prime-power stage \(q>B\), discard the centers from
Lemma~\ref{pairs} conflicting with the initial family.
Its two chains exclude \(O(\log\log q)\) centers in the stage
band; its remaining supports exclude \(O(B^{3/4})\).
These bounds are uniform in their positions and the chain lengths.
Thus every stage still has at least \(9s(q)\) centers when
\(B\) is sufficiently large, since
\[
             B^{3/4}+\log\log q+s(q)=o(q/\log q)
                 \quad\text{uniformly for }q>B.
\]
For any finite cutoff \(X>B\), Hall's theorem assigns \(3s(q)\)
distinct centers to each stage \(B<q\le X\), without sharing centers.
Indeed, a center belongs to at most three stages, so for any
set \(S\) of stages,
\[
             |N(S)|\ge\frac13\sum_{q\in S}\deg(q)
                       \ge\sum_{q\in S}3s(q).
\]
Apply the marriage theorem to \(3s(q)\) copies of each stage.
Distinct centers on the \(8\Z\) grid give separated pairs.
Use \(s(q)\) pairs for the quotient cover and \(2s(q)\) as
the further pairs in \eqref{transfer}.

Starting from \eqref{initial}, the count interval is updated by
\[
                   [u,v]\longmapsto[u+s(q),v+2s(q)].
\]
These updates are always permitted. The next prime power after
\(x\) is at most \(2x\), because a power of two lies in that range;
its \(s\)-value is at most \(2s(x)\).
The initial width \(2s(B)\) therefore suffices for the first stage.
If the old width \(W\) is at least \(s(q)\), the new width is
\(W+s(q)\ge2s(q)\), enough for the next stage.
Successive count intervals overlap and their right endpoints
tend to infinity. Every tail of the sequence consequently covers
all sufficiently large integers. The intervals depend only on
the quotas, although the chosen pools may depend on \(X\).

Each stage-\(q\) pool has total weight \(O(\log q/q^{5/4})\).
Choose \(B\) so large that the sum of these bounds over all
integers \(q>B\) is below \(1/20\), and all preceding uniform
requirements hold. This choice precedes the initial construction:
none of the bounds depends on its later chain lengths or positions.
Every resulting configuration has weight below \(2\).

Finally choose \(X_0>B\) with \(\beta\in G_{X_0}\), possible
because \(\bigcup_nG_n=\Q/\Z\).
For every sufficiently large \(k\), some cutoff \(X\ge X_0\)
has \(k\) in its count interval. The construction at that cutoff
realizes \(G_X\) at count \(k\).
Choose residue zero: its weight is a positive integer below \(2\),
hence \(1\). All intervals longer than two belong to the fixed initial family.
This proves Theorem~\ref{main}.

\section*{AI-declaration}

This version is based on a simplication by GPT6 of~\cite{Land,Tang}.

\begin{thebibliography}{9}
\bibitem{EP289}
T.~F.~Bloom,
\emph{Erd\H{o}s Problem \#289},
\url{https://www.erdosproblems.com/289}.
\bibitem{EG} P.~Erd\H{o}s and R.~L.~Graham,
\emph{Old and New Problems and Results in Combinatorial Number Theory},
Monographies de L'Enseignement Math\'ematique 28, 1980.
\bibitem{HW} G.~H.~Hardy and E.~M.~Wright,
\emph{An Introduction to the Theory of Numbers}, 6th ed.,
Oxford University Press, 2008.
\bibitem{Land} J.~Land,
\emph{A solution of Erd\H{o}s Problem 289 with intervals of length two and three},
supplied manuscript, 4 September 2026.
\bibitem{Tang} Y.~Tang,
\emph{Reciprocal Sums over Separated Integer Intervals},
supplied manuscript, September 2026.
\end{thebibliography}
\end{document}
